% Proposal CODE 6.0 - Software Development
% Aplikasi Manajemen Food Waste

% Document Class
\documentclass[12pt,a4paper,indonesian]{article}

% Packages
\usepackage[utf8]{inputenc}
\usepackage[T1]{fontenc}
\usepackage{mathptmx} % Times New Roman
\usepackage{courier}  % Courier for monospace
\usepackage{helvet}   % Helvetica for sans-serif
\usepackage{amsmath}
\usepackage{amssymb}
\usepackage{pifont}
\usepackage{babel}
\babelprovide[import, main]{indonesian}
\babelprovide[import]{english}
\selectlanguage{indonesian}
\usepackage{geometry}
\usepackage{setspace}
\usepackage{titlesec}
\usepackage{graphicx}
\usepackage{tocloft}
\usepackage{hyperref}
\usepackage{booktabs}
\usepackage{array}
\usepackage[format=hang,font=footnotesize,labelfont=bf]{caption}
\usepackage{enumitem}
\usepackage{parskip}
\usepackage{fancyhdr}
\usepackage{tabularx}
\usepackage{longtable}
\usepackage{ragged2e}
\usepackage{csquotes}
\usepackage{multicol}
\usepackage{multirow}
\usepackage{float}

\newcolumntype{Y}{>{\RaggedRight\arraybackslash}X}
\DeclareQuoteStyle{indonesian}
  {\textquotedblleft}
  {\textquotedblright}
  [0.05em]
  {\textquoteleft}
  {\textquoteright}

% Page Setup - sesuai template: minimal 2.5cm semua sisi
\geometry{
    top=2.5cm,
    bottom=2.5cm,
    left=3cm,
    right=2.5cm
}

% Title Formatting
\titleformat{\section}
    {\normalfont\fontsize{14}{17}\bfseries\centering}{\thesection}{1em}{}

\titleformat{\subsection}
    {\normalfont\fontsize{12}{14}\bfseries}{\thesubsection}{1em}{}

\titleformat{\subsubsection}
    {\normalfont\fontsize{11}{13}\bfseries}{\thesubsubsection}{1em}{}

% Spacing - 1.5 sesuai template
\onehalfspacing

% Header/Footer
\pagestyle{fancy}
\fancyhf{}
\fancyfoot[C]{\thepage}
\renewcommand{\headrulewidth}{0pt}

% Hyperref setup
\hypersetup{
    colorlinks=true,
    linkcolor=black,
    citecolor=blue,
    filecolor=black,
    urlcolor=blue,
}

% Table of Contents customization
\renewcommand{\contentsname}{DAFTAR ISI}
\renewcommand{\cfttoctitlefont}{\hfill\fontsize{14}{17}\selectfont\bfseries}
\renewcommand{\cftaftertoctitle}{\hfill}

% List of Tables customization
\renewcommand{\listtablename}{DAFTAR TABEL}
\renewcommand{\cftlottitlefont}{\hfill\fontsize{14}{17}\selectfont\bfseries}
\renewcommand{\cftafterlottitle}{\hfill}

% List of Figures customization
\renewcommand{\listfigurename}{DAFTAR GAMBAR}
\renewcommand{\cftloftitlefont}{\hfill\fontsize{14}{17}\selectfont\bfseries}
\renewcommand{\cftafterloftitle}{\hfill}

% TOC spacing for Roman numerals (BAB I, II, III...)
\setlength{\cftsecnumwidth}{6em}

% Subsection numbering format
% Format section numbering: BAB I, BAB II, BAB III...
\renewcommand{\thesection}{BAB \Roman{section}}
\renewcommand{\thesubsection}{\arabic{section}.\arabic{subsection}}

% Bibliography Style - IEEE NUMERIC
\usepackage[
    backend=biber,
    style=numeric,
    sorting=none,
    maxbibnames=99,
    hyperref=true,
    backref=false,
    url=false,
    doi=true,
]{biblatex}

\urlstyle{same}

\addbibresource{references.bib}
\addbibresource{include.bib}

% DO NOT redefine \cite - let it use default numeric style

% IEEE style does not need custom delimiters

% Remove "In:" from bibliography
\renewbibmacro{in:}{}

% Ensure pages are not bold
\DeclareFieldFormat{pages}{\normalfont #1}

% Custom format untuk memastikan ada spasi setelah "in"
\DeclareFieldFormat{booktitle}{\addspace\textit{#1}}
\DeclareFieldFormat{journaltitle}{\addspace\textit{#1}}

% Fix URL breaking untuk mencegah overflow
\setcounter{biburllcpenalty}{7000}
\setcounter{biburlucpenalty}{8000}

% Set bibliography to use ragged right alignment
\emergencystretch=1em
\appto{\bibsetup}{\sloppy\RaggedRight}

% Custom Commands
\makeatletter
\newcommand{\myTitle}[1]{\def\@myTitle{#1}}
\newcommand{\myName}[1]{\def\@myName{#1}}
\newcommand{\myNIM}[1]{\def\@myNIM{#1}}
\newcommand{\myTeam}[1]{\def\@myTeam{#1}}
\newcommand{\myUniversity}[1]{\def\@myUniversity{#1}}
\newcommand{\myStudyProgram}[1]{\def\@myStudyProgram{#1}}
\newcommand{\myYear}[1]{\def\@myYear{#1}}

% Initialize default values
\def\@myTitle{ResQfood: Platform Manajemen Food Waste Berbasis Web untuk UMKM Kuliner}
\def\@myTeam{Tim XYZ}
\def\@myNIM{NIM Anggota 1, NIM Anggota 2, NIM Anggota 3, NIM Anggota 4, NIM Anggota 5}
\def\@myName{Nama Anggota 1, Nama Anggota 2, Nama Anggota 3, Nama Anggota 4, Nama Anggota 5}
\def\@myUniversity{Universitas ABC}
\def\@myStudyProgram{Program Studi Teknik Informatika}
\def\@myYear{2026}

\renewcommand{\maketitle}{
    \begin{titlepage}
        \centering
        \vspace*{1cm}
        
        % Judul Aplikasi
        {\fontsize{16}{20}\selectfont
        \textbf{RESQFOOD}\par}
        
        \vspace{0.5cm}
        
        {\fontsize{12}{14}\selectfont
        \textbf{Platform Manajemen Food Waste Berbasis Web}\\[0.3cm]
        \textbf{untuk UMKM Kuliner}\par}
        
        \vspace{2cm}
        
        % Nama Tim
        {\fontsize{12}{14}\selectfont
        \textbf{Kelompok: \@myTeam}\par}
        
        \vspace{1cm}
        
        % Anggota Tim
        {\fontsize{11}{13}\selectfont
        \@myName\par}
        
        \vspace{0.5cm}
        
        {\fontsize{11}{13}\selectfont
        \@myNIM\par}
        
        \vspace{2cm}
        
        % Logo
        \IfFileExists{logo.png}{
            \includegraphics[width=0.3\textwidth]{logo.png}
        }{
            % Placeholder jika logo belum ada
        }
        
        \vfill
        
        % Informasi Universitas
        {\fontsize{12}{14}\selectfont
        \textbf{\@myStudyProgram}\par
        \textbf{\@myUniversity}\par}
        
        \vspace{0.5cm}
        
        {\fontsize{14}{17}\selectfont
        \textbf{\@myYear}\par}
        
    \end{titlepage}
    \newpage
}
\makeatother

% Beginning of Document
\begin{document}
\hypersetup{pageanchor=false}

% Force TOC title to be uppercase
\renewcommand{\contentsname}{DAFTAR ISI}
\renewcommand{\listtablename}{DAFTAR TABEL}
\renewcommand{\listfigurename}{DAFTAR GAMBAR}

% Use Times New Roman as default font
\rmfamily

% Title Page
\maketitle

% Halaman Depan menggunakan angka Romawi
\pagenumbering{roman}

% DAFTAR ISI
\phantomsection
\addcontentsline{toc}{section}{DAFTAR ISI}
\tableofcontents
\clearpage

% DAFTAR TABEL
\phantomsection
\addcontentsline{toc}{section}{DAFTAR TABEL}
\listoftables
\clearpage

% DAFTAR GAMBAR
\phantomsection
\addcontentsline{toc}{section}{DAFTAR GAMBAR}
\listoffigures
\clearpage

% ==================== BAGIAN 1: COVER ====================
% (Cover sudah ditangani oleh \maketitle di atas)

% ==================== BAGIAN 2: ABSTRAK ====================
\clearpage
\setcounter{section}{0}
\pagenumbering{arabic}
\setcounter{page}{1}
\hypersetup{pageanchor=true}

\section*{ABSTRAK}
\phantomsection
\addcontentsline{toc}{section}{ABSTRAK}
Indonesia menghasilkan sekitar 23-48 juta ton sampah makanan per tahun, dengan proporsi terbesar berasal dari sektor UMKM kuliner yang kehilangan 10-30\% produk akibat surplus yang tidak terjual. Masalah ini tidak hanya menimbulkan kerugian ekonomi bagi pelaku UMKM, tetapi juga berkontribusi signifikan terhadap emisi gas rumah kaca dan pemborosan sumber daya alam. Sayangnya, belum ada solusi terpadu yang secara efektif menjembatani UMKM kuliner dengan konsumen dan mitra donasi untuk menyalurkan makanan surplus yang masih layak konsumsi.

ResQfood hadir sebagai platform manajemen \textit{food waste} berbasis web yang menghubungkan UMKM kuliner dengan konsumen dan mitra donasi. Platform ini menawarkan serangkaian fitur inovatif meliputi \textit{Food Eligibility Assessment} untuk menilai kelayakan makanan secara objektif, \textit{Dynamic Pricing} untuk mengoptimalkan harga jual berdasarkan kondisi makanan, \textit{Sistem Prediksi Produksi Harian} berbasis data historis dan cuaca, serta mekanisme donasi otomatis melalui \textit{Food Rescue Score}. Dengan memanfaatkan teknologi \textit{machine learning} untuk prediksi produksi dan analitik \textit{real-time}, ResQfood bertujuan mengurangi \textit{food waste} hingga 40\% sekaligus meningkatkan pendapatan UMKM melalui penjualan makanan surplus.

Platform ini memberikan dampak positif bagi tiga pihak utama: UMKM mendapatkan pendapatan tambahan dari makanan yang sebelumnya terbuang, konsumen memperoleh akses terhadap makanan berkualitas dengan harga terjangkau, serta lingkungan mendapat manfaat dari pengurangan emisi gas rumah kaca dan volume sampah organik. Dengan target implementasi di Kota Bandung pada tahun pertama, ResQfood berpotensi menjadi solusi \textit{scalable} yang dapat diterapkan di kota-kota besar lainnya di Indonesia.

\clearpage

% ==================== BAGIAN 3: LATAR BELAKANG & URGENSI ====================
% REMOVED: Section header "LATAR BELAKANG & URGENSI MASALAH" 
% (redundant - content already in BAB I PENDAHULUAN)

\subsection{Konteks Masalah}

Indonesia merupakan salah satu negara penghasil sampah makanan terbesar di dunia. Berdasarkan data dari Badan Pusat Statistik (BPS) tahun 2024, Indonesia menghasilkan sekitar 23-48 juta ton sampah makanan per tahun, yang setara dengan 115-237 kg per kapita. Angka ini menjadikan Indonesia peringkat kedua di dunia setelah Arab Saudi dalam hal pemborosan makanan per kapita.

Sektor UMKM kuliner menjadi salah satu kontributor terbesar terhadap \textit{food waste}. Data Kementerian Koperasi dan Usaha Kecil dan Menengah menunjukkan bahwa UMKM kuliner di Indonesia mencapai lebih dari 5 juta unit usaha, dengan tingkat kerugian akibat makanan surplus yang tidak terjual berkisar antara 10-30\% dari total produksi harian. Kerugian ini tidak hanya berdampak pada aspek finansial UMKM, tetapi juga menciptakan masalah lingkungan yang serius.

\begin{table}[H]
\centering
\caption{Data Sampah Makanan di Indonesia}
\label{tab:data_sampah}
\begin{tabularx}{\textwidth}{lX}
\toprule
\textbf{Indikator} & \textbf{Data} \\
\midrule
Total sampah makanan per tahun & 23-48 juta ton \\
Sampah per kapita per tahun & 115-237 kg \\
Kontribusi UMKM kuliner & 10-30\% dari produksi \\
Persentase sampah organik & 50-70\% total sampah \\
Emisi CO$_2$ dari food waste & $\sim$50 juta ton per tahun \\
\bottomrule
\end{tabularx}
\end{table}

\subsection{Akar Masalah}

Berdasarkan analisis yang telah dilakukan, terdapat beberapa akar masalah utama yang menyebabkan tingginya \textit{food waste} di sektor UMKM kuliner:

\begin{enumerate}[leftmargin=*]
    \item \textbf{Ketidakpastian Permintaan} — UMKM kuliner sering kali kesulitan memprediksi jumlah permintaan konsumen dari hari ke hari, sehingga cenderung memproduksi makanan dalam jumlah berlebih untuk mengantisipasi lonjakan permintaan.
    
    \item \textbf{Keterbatasan Infrastruktur Penyimpanan} — Banyak UMKM kuliner yang tidak memiliki fasilitas penyimpanan yang memadai untuk menjaga kesegaran makanan, sehingga makanan yang tidak terjual harus segera dibuang.
    
    \item \textbf{Kurangnya Akses ke Pasar Alternatif} — Tidak adanya platform yang menghubungkan UMKM dengan konsumen potensial untuk menjual makanan surplus dengan harga diskon menjadi hambatan utama.
    
    \item \textbf{Minimnya Kesadaran Lingkungan} — Banyak pelaku UMKM yang belum menyadari dampak lingkungan dari pembuangan makanan berlebih, sehingga kurang termotivasi untuk mengurangi \textit{food waste}.
\end{enumerate}

\subsection{Dampak jika Tidak Diselesaikan}

Jika masalah \textit{food waste} di sektor UMKM kuliner tidak segera ditangani, konsekuensi yang akan terjadi dalam 1-2 tahun ke depan meliputi:

\begin{itemize}[leftmargin=*]
    \item \textbf{Kerugian Ekonomi yang Meningkat} — UMKM akan terus kehilangan pendapatan akibat makanan yang terbuang, yang diperkirakan mencapai Rp 500 ribu hingga Rp 2 juta per hari untuk satu unit UMKM kuliner.
    
    \item \textbf{Peningkatan Emisi Gas Rumah Kaca} — Sampah makanan yang terurai di Tempat Pembuangan Akhir (TPA) menghasilkan metana, gas rumah kaca yang 25 kali lebih potensial dari CO$_2$ dalam mempercepat perubahan iklim.
    
    \item \textbf{Beban Lingkungan yang Semakin Berat} — Volume sampah organik yang terus meningkat akan mempercepat kerusakan TPA dan meningkatkan risiko pencemaran tanah serta air tanah.
\end{itemize}

\subsection{Mengapa Sekarang?}

Beberapa faktor yang membuat penyelesaian masalah ini menjadi sangat mendesak saat ini:

\begin{enumerate}[leftmargin=*]
    \item \textbf{Kebijakan Pemerintah} — Pemerintah Indonesia telah menetapkan target pengurangan sampah makanan sebesar 50\% pada tahun 2030 melalui Peraturan Presiden Nomor 97 Tahun 2017 tentang Kebijakan dan Strategi Nasional Pengelolaan Sampah Rumah Tangga dan Sampah Sejenis Sampah Rumah Tangga.
    
    \item \textbf{Kesadaran Konsumen yang Meningkat} — Generasi milenial dan Gen Z menunjukkan peningkatan kesadaran terhadap isu lingkungan, termasuk \textit{food waste}, sehingga membuka peluang pasar untuk solusi yang ramah lingkungan.
    
    \item \textbf{Perkembangan Teknologi} — Kemajuan teknologi \textit{web} dan \textit{mobile application} memungkinkan pengembangan platform yang mudah diakses oleh UMKM dengan biaya yang terjangkau.
    
    \item \textbf{Dampak Pandemi} — Pandemi COVID-19 telah mengubah pola konsumsi masyarakat, di mana banyak UMKM kuliner yang mengalami fluktuasi permintaan yang signifikan, sehingga kebutuhan akan sistem manajemen \textit{food waste} menjadi semakin mendesak.
\end{enumerate}

\clearpage

% ==================== BAGIAN 4: TUJUAN, MANFAAT, & DAMPAK ====================
\subsection{Tujuan}

ResQfood dirancang dengan tujuan yang spesifik, terukur, dapat dicapai, relevan, dan berbatas waktu (SMART) sebagai berikut:

\begin{enumerate}[leftmargin=*]
    \item \textbf{Specific (Spesifik)} — Mengembangkan platform manajemen \textit{food waste} berbasis web yang menghubungkan UMKM kuliner dengan konsumen dan mitra donasi untuk menyalurkan makanan surplus yang masih layak konsumsi.
    
    \item \textbf{Measurable (Terukur)} — Mengurangi volume \textit{food waste} pada UMKM kuliner mitra sebesar 40\% dalam 6 bulan pertama operasional, serta meningkatkan pendapatan UMKM dari penjualan makanan surplus sebesar 25\%.
    
    \item \textbf{Achievable (Dapat Dicapai)} — Menerapkan fitur \textit{Food Eligibility Assessment}, \textit{Dynamic Pricing}, dan \textit{Sistem Prediksi Produksi} yang telah terbukti efektif dalam studi literatur dan pilot project.
    
    \item \textbf{Relevant (Relevan)} — Menjawab kebutuhan mendesak akan solusi \textit{food waste} yang terintegrasi, sesuai dengan target pengurangan sampah makanan nasional sebesar 50\% pada tahun 2030.
    
    \item \textbf{Time-bound (Berbatas Waktu)} — Mencapai target implementasi di 100 UMKM kuliner di Kota Bandung dalam 12 bulan pertama, dengan ekspansi ke kota lain pada tahun kedua.
\end{enumerate}

\subsection{Manfaat}

Manfaat ResQfood dapat dilihat dari tiga perspektif utama:

\subsubsection{Manfaat bagi UMKM Kuliner}
\begin{itemize}[leftmargin=*]
    \item Mengurangi kerugian akibat makanan surplus yang tidak terjual
    \item Mendapatkan pendapatan tambahan melalui penjualan makanan surplus dengan harga diskon
    \item Meningkatkan efisiensi operasional melalui prediksi produksi yang akurat
    \item Meningkatkan reputasi sebagai UMKM yang peduli lingkungan
    \item Mendapatkan akses ke jaringan donasi untuk makanan yang tidak terjual
\end{itemize}

\subsubsection{Manfaat bagi Konsumen}
\begin{itemize}[leftmargin=*]
    \item Memperoleh akses terhadap makanan berkualitas dengan harga yang lebih terjangkau
    \item Berkontribusi langsung dalam pengurangan \textit{food waste} dan emisi gas rumah kaca
    \item Mendapatkan informasi transparan mengenai kelayakan dan asal-usul makanan
    \item Memperoleh penghargaan berupa reward dan badge atas kontribusi lingkungan
\end{itemize}

\subsubsection{Manfaat bagi Lingkungan}
\begin{itemize}[leftmargin=*]
    \item Mengurangi volume sampah organik yang berakhir di TPA
    \item Menurunkan emisi gas rumah kaca dari pembusukan sampah makanan
    \item Melestarikan sumber daya alam yang digunakan dalam produksi makanan
    \item Mendorong terciptanya budaya pengurangan \textit{food waste} di masyarakat
\end{itemize}

\subsection{Dampak Jangka Panjang}

\subsubsection{Dampak Ekonomi}
Dalam jangka panjang, ResQfood berpotensi menciptakan ekosistem ekonomi sirkular di sektor kuliner, di mana setiap produk makanan memiliki nilai ekonomis hingga akhir siklus hidupnya. Hal ini tidak hanya menguntungkan UMKM dan konsumen, tetapi juga membuka peluang kerja baru di bidang logistik dan manajemen \textit{food waste}.

\subsubsection{Dampak Sosial}
Platform ini berpotensi memperkuat hubungan antara UMKM dengan komunitas sekitar melalui mekanisme donasi. Makanan surplus yang didonasikan dapat membantu memenuhi kebutuhan pangan kelompok masyarakat yang membutuhkan, sehingga berkontribusi pada penurunan angka ketahanan pangan.

\subsubsection{Dampak Lingkungan}
Dengan target pengurangan \textit{food waste} sebesar 40\%, ResQfood berkontribusi pada pengurangan emisi gas rumah kaca dan pelestarian sumber daya alam. Skalabilitas platform ini memungkinkan dampak lingkungan yang lebih luas saat diterapkan di berbagai kota di Indonesia.

\subsection{Metrik Dampak}

\begin{table}[H]
\centering
\caption{Metrik Dampak ResQfood}
\label{tab:metrik_dampak}
\begin{tabularx}{\textwidth}{lXX}
\toprule
\textbf{Metrik} & \textbf{Target 6 Bulan} & \textbf{Target 12 Bulan} \\
\midrule
Jumlah UMKM mitra & 50 UMKM & 100 UMKM \\
Volume food waste terkurangi & 20\% & 40\% \\
Pendapatan tambahan UMKM & 15\% & 25\% \\
Jumlah makanan diselamatkan & 5.000 kg & 15.000 kg \\
Pengurangan emisi CO$_2$ & 10 ton & 30 ton \\
Jumlah konsumen aktif & 1.000 & 5.000 \\
Jumlah donasi tersalurkan & 1.000 kg & 5.000 kg \\
\bottomrule
\end{tabularx}
\end{table}

\clearpage

% ==================== BAGIAN 5: NILAI INOVASI & ORISINALITAS ====================
\section{DIFERENSIASI DAN ANALISIS PASAR}

\subsection{Diferensiasi Solusi}

Savora memiliki beberapa keunggulan diferensiasi yang membedakannya dari solusi \textit{food waste} yang sudah ada:

\begin{enumerate}[leftmargin=*]
    \item \textbf{Pendekatan Terpadu} — Berbeda dengan platform yang hanya fokus pada satu aspek (misalnya hanya penjualan diskon atau hanya donasi), Savora mengintegrasikan seluruh rantai nilai \textit{food waste} mulai dari penilaian kelayakan makanan, penjualan surplus dengan harga dinamis, hingga pelacakan dampak lingkungan \cite{kusumawati2025mobile}\cite{mobilefoodwaste2024review}.
    
    \item \textbf{Food Eligibility Assessment} — Sistem penilaian kelayakan makanan yang objektif dan transparan, memberikan kepastian bagi konsumen mengenai kualitas dan keamanan makanan yang dibeli. Menggunakan \textit{weighted scoring algorithm} berbasis 4 parameter utama.
    
    \item \textbf{Dynamic Pricing Engine} — Mesin penetapan harga dinamis yang mengoptimalkan harga berdasarkan sisa waktu kelayakan, permintaan pasar, dan margin UMKM \cite{qlearning2026pricing}\cite{kayikci2024iot}, memaksimalkan peluang penjualan sebelum makanan kadaluarsa.
    
    \item \textbf{Transparansi \& Trust Building} — Sistem \textit{Food Trust Index} yang menampilkan rating dan ulasan konsumen terhadap UMKM, serta \textit{Carbon Impact Tracking} yang menunjukkan dampak positif setiap pembelian terhadap lingkungan.
\end{enumerate}

\subsection{Keunggulan Teknis}

\subsubsection{Food Eligibility Score}
Sistem ini menggunakan model \textit{scoring} berbasis aturan yang mempertimbangkan:

\begin{itemize}[leftmargin=*]
    \item Waktu produksi dan estimasi masa simpan
    \item Kondisi kemasan dan suhu penyimpanan
    \item Riwayat penanganan makanan
    \item Kategori makanan dan karakteristik spesifik
\end{itemize}

\subsubsection{Dynamic Pricing Engine}
Mesin penetapan harga dinamis yang mengoptimalkan:

\begin{itemize}[leftmargin=*]
    \item Harga berdasarkan sisa waktu kelayakan
    \item Elasticitas permintaan lokal
    \item Harga kompetitor di area yang sama
    \item Margin minimum yang menguntungkan UMKM
\end{itemize}

\subsection{Technology Stack yang Unik}

Kombinasi teknologi yang digunakan dalam Savora memberikan keunggulan kompetitif:

\begin{table}[H]
\centering
\caption{Technology Stack Savora}
\label{tab:tech_stack}
\begin{tabularx}{\textwidth}{|l|Y|Y|}
\hline
\textbf{Komponen} & \textbf{Teknologi} & \textbf{Alasan Pemilihan} \\
\hline
Frontend & React.js & Komponen reusable, performa tinggi \\
\hline
Backend & Node.js + Express & asynchronous, cocok untuk real-time \\
\hline
Database & PostgreSQL + Redis & ACID compliance + caching \\
\hline
ML & Python + scikit-learn & Ekosistem ML yang mature \\
\hline
Real-time & Socket.IO & Komunikasi dua arah \\
\hline
Hosting & AWS / GCP & Skalabilitas dan reliability \\
\hline
\end{tabularx}
\end{table}

\subsection{Moat (Penghalang Kompetisi)}

Savora memiliki beberapa \textit{moat} yang melindungi posisi kompetitif:

\begin{enumerate}[leftmargin=*]
    \item \textbf{Data Lokal} — Dataset penjualan UMKM kuliner Indonesia yang dikumpulkan dari mitra merupakan aset berharga yang sulit ditiru oleh kompetitor baru atau perusahaan besar.
    
    \item \textbf{Jaringan UMKM} — Hubungan langsung dengan UMKM kuliner lokal memberikan keunggulan dalam hal \textit{network effect} dan \textit{switching cost}.
    
    \item \textbf{Algoritma yang Dioptimasi} — Model \textit{machine learning} yang dilatih dengan data lokal memiliki akurasi yang lebih tinggi dibandingkan model generik.
    
    \item \textbf{Mekanisme Donasi} — Integrasi dengan jaringan donasi memberikan nilai tambah yang sulit ditiru oleh platform yang hanya fokus pada penjualan.
\end{enumerate}

\clearpage

% ==================== BAGIAN 6: ANALISIS KOMPETITOR & BUSINESS VIABILITY ====================
\section{ANALISIS KOMPETITOR}

\subsection{Analisis Kompetitor}

Berikut adalah perbandingan Savora dengan platform serupa yang sudah ada di pasar berdasarkan data aktual per 2025--2026 \cite{mobilefoodwaste2024review}\cite{kusumawati2025mobile}:

\begin{landscape}
\begin{table}[H]
\centering
\caption{Perbandingan Fitur dengan Kompetitor}
\label{tab:kompetitor}
\small
\begin{tabularx}{\linewidth}{lXXXXX}
\toprule
\textbf{Aspek / Fitur}
  & \textbf{Savora}
  & \textbf{Too Good To Go}
  & \textbf{Surplus Indonesia}
  & \textbf{Garda Pangan}
  & \textbf{Olio} \\
\midrule
Jenis platform
  & Marketplace food rescue UMKM (MVP)
  & Global marketplace food rescue
  & Clearance sale app (recommerce)
  & Food bank sosial / nirlaba
  & P2P sharing + solusi bisnis \\
\addlinespace
Tersedia di Indonesia
  & \checkmark\ (MVP)
  & \ding{55}\ (baru Jepang, 2026)
  & \checkmark\ (aktif)
  & \checkmark\ (Surabaya \& sekitar)
  & Terbatas (komunitas minimal) \\
\addlinespace
Fokus UMKM lokal kuliner
  & \checkmark
  & Terbatas (resto besar)
  & \checkmark\ (UMKM, horeka)
  & \checkmark\ (distribusi gratis)
  & \ding{55} \\
\addlinespace
Transparansi produk
  & Tinggi (foto, deskripsi, kelayakan)
  & Rendah (\textit{Surprise Bag})
  & Sedang (nama \& foto)
  & Tinggi (SOP seleksi ketat)
  & Sedang \\
\addlinespace
Food Trust Index / info kelayakan
  & \checkmark\ (rule-based, tampil ke user)
  & \ding{55}
  & \ding{55}
  & Seleksi internal (tidak ke user)
  & \ding{55} \\
\addlinespace
Metode pembayaran
  & COD + payment gateway (Midtrans)
  & Digital in-app (tanpa COD)
  & Digital saja (OVO, GoPay, DANA; tanpa COD)
  & Tidak ada (donasi)
  & Digital in-app \\
\addlinespace
Dynamic discount berbasis kelayakan
  & \checkmark\ (rule-based, 10--50\%)
  & \checkmark\ (AI, berbasis historis, sejak 2025)
  & \checkmark\ (hingga 80\%, tidak berbasis kelayakan)
  & \ding{55}
  & \ding{55} \\
\addlinespace
Pickup code \& batas waktu
  & \checkmark\ (pickup code + deadline)
  & Window waktu (tanpa pickup code)
  & Window waktu
  & \ding{55}
  & \ding{55} \\
\addlinespace
Waste Log untuk mitra UMKM
  & \checkmark
  & \ding{55}
  & \ding{55}
  & \ding{55}
  & \ding{55} \\
\addlinespace
Help Center food rescue
  & \checkmark\ (laporan terstruktur)
  & Ada (umum)
  & Tidak dipublikasikan khusus
  & \ding{55}
  & \ding{55} \\
\addlinespace
Impact tracking ke pengguna
  & \checkmark\ (estimatif, bukan klaim resmi)
  & \checkmark\ (CO\textsubscript{2}e per transaksi)
  & \checkmark\ (estimasi CO\textsubscript{2}/CH\textsubscript{4})
  & \ding{55}
  & \checkmark\ (3,67 kg CO\textsubscript{2}e per kg) \\
\addlinespace
Rating \& ulasan
  & \checkmark
  & \checkmark
  & \checkmark
  & \ding{55}
  & \checkmark \\
\addlinespace
Sertifikasi
  & ---
  & B-Corp
  & B-Corp (pertama di Indonesia)
  & --- (nirlaba)
  & B-Corp \\
\addlinespace
Skala (2025--2026)
  & MVP lomba
  & 120 juta user, 20+ negara
  & Ratusan ribu user, Jawa/Bali/Sulsel
  & Lokal Surabaya
  & 9 juta user, 62+ negara \\
\bottomrule
\end{tabularx}
\end{table}
\end{landscape}

\subsection{Analisis SWOT}

\begin{table}[H]
\centering
\caption{Analisis SWOT Savora}
\label{tab:swot}
\begin{tabularx}{\textwidth}{XX}
\toprule
\textbf{Strengths (Kekuatan)} & \textbf{Weaknesses (Kelemahan)} \\
\midrule
\begin{itemize}[leftmargin=*,nosep]
    \item Fitur lengkap terintegrasi (FES, Dynamic Pricing, Trust Index)
    \item Fokus pada UMKM lokal Indonesia
    \item Transparansi kelayakan makanan
    \item Mekanisme donasi terintegrasi
\end{itemize} \\ 
& \begin{itemize}[leftmargin=*,nosep]
    \item Butuh waktu untuk adopsi awal
    \item Bergantung pada jumlah UMKM mitra
    \item Membutuhkan data historis yang cukup
\end{itemize} \\
\addlinespace
\textbf{Opportunities (Peluang)} & \textbf{Threats (Ancaman)} \\
\midrule
\begin{itemize}[leftmargin=*,nosep]
    \item Kebijakan pengurangan food waste nasional
    \item Peningkatan kesadaran lingkungan
    \item Potensi ekspansi ke kota lain
    \item Kolaborasi dengan pemerintah daerah
\end{itemize} \\ 
& \begin{itemize}[leftmargin=*,nosep]
    \item Masuknya pemain besar dari luar negeri
    \item Perubahan regulasi
    \item Resistensi UMKM terhadap teknologi baru
\end{itemize} \\
\bottomrule
\end{tabularx}
\end{table}

\subsection{Model Bisnis}

Savora mengadopsi model bisnis \textit{freemium} dengan beberapa sumber pendapatan:

\subsubsection{Revenue Streams}
\begin{enumerate}[leftmargin=*]
    \item \textbf{Komisi Transaksi} — Sebesar 5-10\% dari setiap transaksi penjualan makanan surplus melalui platform.
    \item \textbf{Langganan Premium} — Paket premium untuk UMKM dengan fitur analitik lanjutan, prediksi produksi, dan laporan detail seharga Rp 99.000/bulan.
    \item \textbf{Iklan dan Promosi} — Layanan promosi produk untuk UMKM mitra yang ingin meningkatkan visibilitas.
    \item \textbf{Kerjasama Korporat} — Program CSR perusahaan yang ingin berkontribusi pada pengurangan \textit{food waste}.
\end{enumerate}

\subsubsection{Pricing Strategy}
\begin{table}[H]
\centering
\caption{Strategi Harga Savora}
\label{tab:pricing}
\begin{tabularx}{\textwidth}{lXX}
\toprule
\textbf{Paket} & \textbf{Harga} & \textbf{Fitur} \\
\midrule
Basic & Gratis & Fitur dasar, maks 50 produk/bulan \\
Pro & Rp 99.000/bulan & Analitik lanjutan, laporan detail, unlimited produk \\
Enterprise & Custom & Khusus jaringan UMKM besar, dedicated support \\
\bottomrule
\end{tabularx}
\end{table}

\subsection{Rencana Sustainability}

\begin{enumerate}[leftmargin=*]
    \item \textbf{Tahun Pertama} — Fokus pada akuisisi UMKM dan pengguna di Kota Bandung. Target mencapai \textit{break even} dari komisi transaksi dalam 12 bulan.
    
    \item \textbf{Tahun Kedua} — Ekspansi ke 3-5 kota besar lainnya (Jakarta, Surabaya, Yogyakarta). Meluncurkan paket premium dan program CSR.
    
    \item \textbf{Tahun Ketiga} — Menjadi platform referensi manajemen \textit{food waste} di Indonesia. Potensi pendanaan dari investor untuk percepatan pertumbuhan.
    
    \item \textbf{Jangka Panjang} — Ekspansi ke Asia Tenggara dan pengembangan produk turunan seperti pelatihan UMKM dan konsultasi \textit{sustainability}.
\end{enumerate}

\clearpage

% ==================== BAGIAN 7: TEKNOLOGI YANG DIGUNAKAN ====================
\section{TEKNOLOGI}

\subsection{Daftar Teknologi}

Berikut adalah daftar lengkap teknologi yang digunakan dalam pengembangan Savora:

\begin{table}[H]
\centering
\caption{Daftar Teknologi Savora}
\label{tab:daftar_teknologi}
\small
\begin{tabularx}{\textwidth}{|l|X|X|}
\hline
\textbf{Kategori} & \textbf{Teknologi} & \textbf{Versi} \\
\hline
Bahasa Pemrograman & JavaScript, TypeScript & ES2024, 5.x \\
\hline
Frontend Framework & React.js & 18.x \\
\hline
Backend Framework & Node.js + Express & 20 LTS, 4.x \\
\hline
Database & PostgreSQL, Redis & 16.x, 7.x \\
\hline
ORM & Prisma & 5.x \\
\hline
Authentication & JWT, bcrypt & - \\
\hline
Real-time & Socket.IO & 4.x \\
\hline
Machine Learning & Python, scikit-learn & 3.12, 1.4 \\
\hline
API Design & RESTful API & - \\
\hline
Version Control & Git, GitHub & - \\
\hline
CI/CD & GitHub Actions & - \\
\hline
Hosting & AWS / Vercel & - \\
\hline
\end{tabularx}
\end{table}

\subsection{Alasan Pemilihan Teknologi}

\subsubsection{React.js untuk Frontend}
React.js dipilih sebagai \textit{frontend framework} karena:
\begin{itemize}[leftmargin=*]
    \item \textbf{Komponen Reusable} — Arsitektur berbasis komponen memungkinkan pembuatan UI yang konsisten dan mudah dipelihara.
    \item \textbf{Performa Tinggi} — Virtual DOM memberikan performa yang optimal untuk aplikasi dengan banyak interaksi.
    \item \textbf{Ekosistem Kaya} — Tersedia banyak library pendukung seperti React Router, Redux, dan Material-UI.
    \item \textbf{Komunitas Besar} — Mudah menemukan dokumentasi, tutorial, dan solusi untuk berbagai masalah.
\end{itemize}

\subsubsection{Node.js + Express untuk Backend}
Node.js dan Express dipilih karena:
\begin{itemize}[leftmargin=*]
    \item \textbf{Asynchronous} — Sangat cocok untuk aplikasi yang membutuhkan pemrosesan banyak request secara simultan.
    \item \textbf{JavaScript Full-stack} — Penggunaan bahasa yang sama di frontend dan backend mempercepat pengembangan.
    \item \textbf{Ringan dan Cepat} — Cocok untuk microservice dan API yang responsif.
    \item \textbf{NPM Ecosystem} — Akses ke ribuan package yang sudah teruji.
\end{itemize}

\subsubsection{PostgreSQL + Redis}
\begin{itemize}[leftmargin=*]
    \item \textbf{PostgreSQL} — Database relasional yang robust, mendukung ACID, dan memiliki fitur JSON untuk data semi-terstruktur.
    \item \textbf{Redis} — In-memory database untuk caching, session management, dan real-time features.
\end{itemize}

\subsubsection{Python + scikit-learn untuk Machine Learning}
Python dipilih untuk komponen \textit{machine learning} karena:
\begin{itemize}[leftmargin=*]
    \item \textbf{Ekosistem ML Terkaya} — scikit-learn, pandas, dan numpy menyediakan tools yang komprehensif untuk algoritma scoring dan analitik.
    \item \textbf{Integrasi Mudah} — Dapat diintegrasikan dengan backend Node.js melalui API atau \textit{microservice}.
    \item \textbf{Prototyping Cepat} — Memungkinkan eksperimen dan iterasi model dengan cepat.
\end{itemize}

\subsection{Arsitektur Teknologi}

\begin{figure}[H]
\centering
% Placeholder untuk diagram arsitektur
\fbox{\parbox{0.8\textwidth}{\centering\vspace{3cm}Diagram Arsitektur Sistem\\(Akan diganti dengan diagram aktual)\vspace{3cm}}}
\caption{Arsitektur High-Level Savora}
\label{fig:arsitektur}
\end{figure}

Alur komunikasi sistem:
\begin{enumerate}[leftmargin=*]
    \item \textbf{Client Layer} — Browser pengguna mengakses React.js application.
    \item \textbf{API Gateway} — Request diteruskan ke Express.js backend.
    \item \textbf{Application Layer} — Backend memproses business logic, autentikasi, dan validasi data.
    \item \textbf{Data Layer} — PostgreSQL menyimpan data persisten, Redis menyimpan cache dan session.
    \item \textbf{ML Service} — Python microservice menangani Food Eligibility Score calculation dan analitik dashboard.
    \item \textbf{Real-time Layer} — Socket.IO menangani notifikasi dan update status secara real-time.
\end{enumerate}

\clearpage

% ==================== BAGIAN 8: BATASAN PERANGKAT LUNAK ====================
\section{BATASAN PERANGKAT LUNAK}

Bagian ini akan dilengkapi setelah aplikasi selesai dikembangkan dan siap dikumpulkan.

\clearpage

% ==================== BAGIAN 9: METODOLOGI PENGEMBANGAN ====================
\section{METODOLOGI PENGEMBANGAN}

\subsection{Metodologi yang Dipilih}

Savora dikembangkan menggunakan metodologi \textbf{Scrum} dengan iterasi sprint 2 minggu. Metodologi ini dipilih karena fleksibilitasnya dalam mengakomodasi perubahan kebutuhan dan kemampuannya untuk menghasilkan produk yang selaras dengan ekspektasi pengguna.

\subsection{Alasan Pemilihan Scrum}

\begin{enumerate}[leftmargin=*]
    \item \textbf{Iterasi Cepat} : Memungkinkan pengembangan dan pengujian fitur secara bertahap, sehingga bug dapat teridentifikasi lebih awal.
    \item \textbf{Feedback Berkala} : Stakeholder dapat memberikan masukan setiap akhir sprint untuk perbaikan berkelanjutan.
    \item \textbf{Transparansi} : Daily standup dan sprint review memastikan seluruh tim memiliki pemahaman yang sama tentang progres.
    \item \textbf{Adaptif} : Perubahan prioritas atau kebutuhan dapat diakomodasi tanpa mengganggu seluruh roadmap.
\end{enumerate}

\subsection{Tahapan dan Timeline Pengembangan}

\begin{table}[H]
\centering
\caption{Tahapan Pengembangan Savora}
\label{tab:tahapan}
\begin{tabularx}{\textwidth}{|l|X|X|}
\hline
\textbf{Fase} & \textbf{Aktivitas} & \textbf{Durasi} \\
\hline
Sprint 0 & Setup environment, Wireframe, Database design & 1 minggu \\
\hline
Sprint 1 & Autentikasi, Dashboard Admin & 2 minggu \\
\hline
Sprint 2 & Manajemen Produk UMKM, Marketplace & 2 minggu \\
\hline
Sprint 3 & Food Eligibility, Dynamic Pricing & 2 minggu \\
\hline
Sprint 4 & Checkout, Payment simulation & 2 minggu \\
\hline
Sprint 5 & Tracking, Rating, Notifikasi & 2 minggu \\
\hline
Sprint 6 & Integrasi, Testing, Deployment & 2 minggu \\
\hline
\end{tabularx}
\end{table}

\begin{figure}[H]
\centering
% Placeholder untuk Gantt chart
\fbox{\parbox{0.8\textwidth}{\centering\vspace{3cm}Gantt Chart Pengembangan\\(Akan diganti dengan Gantt chart aktual)\vspace{3cm}}}
\caption{Timeline Pengembangan Savora}
\label{fig:gantt}
\end{figure}

\subsection{Peran Anggota Tim}

\begin{table}[H]
\centering
\caption{Pembagian Tugas Anggota Tim}
\label{tab:peran}
\begin{tabularx}{\textwidth}{|l|Y|Y|}
\hline
\textbf{Role} & \textbf{Anggota} & \textbf{Responsibilitas} \\
\hline
Project Manager & Wa Ode Nur Alia & Koordinasi, perencanaan, komunikasi \\
\hline
Frontend Developer & Richard Firmansyah & Pengembangan UI/UX, React.js \\
\hline
Backend Developer & Muhammad Rifaidi & API, Database, Autentikasi \\
\hline
ML Engineer & Ridwan Hakim Ramadhan & Food Eligibility Score, Dynamic Pricing, Analitik \\
\hline
QA \& DevOps & Nadi Azzada Akbar & Testing, Deployment, CI/CD \\
\hline
\end{tabularx}
\end{table}

\subsection{Tools Pendukung}

\begin{itemize}[leftmargin=*]
    \item \textbf{Project Management} : Notion atau Trello untuk sprint board dan task tracking
    \item \textbf{Communication} : Discord atau WhatsApp Group untuk komunikasi harian
    \item \textbf{Design} : Figma untuk wireframe dan UI design
    \item \textbf{Documentation} : Google Docs untuk dokumentasi teknis
    \item \textbf{Version Control} : Git dengan branch strategy (main, develop, feature branches)
\end{itemize}

\clearpage

% ==================== BAGIAN 10: ARSITEKTUR SISTEM & DESAIN ====================
\section{ARSITEKTUR SISTEM}

\subsection{Diagram Arsitektur Sistem}

\begin{figure}[H]
\centering
% Placeholder untuk diagram
\fbox{\parbox{0.8\textwidth}{\centering\vspace{4cm}Diagram Use Case\\(Akan diganti dengan diagram aktual)\vspace{4cm}}}
\caption{Use Case Diagram Savora}
\label{fig:usecase}
\end{figure}

Terdapat tiga aktor utama dalam sistem Savora:
\begin{enumerate}[leftmargin=*]
    \item \textbf{Customer} — Pengguna yang membeli makanan surplus, memberikan rating, dan melacak pesanan.
    \item \textbf{UMKM} — Pengelola produk yang mengelola stok, menentukan harga, dan memproses pesanan.
    \item \textbf{Admin} — Pengelola sistem yang memantau aktivitas, mengelola user, dan memverifikasi donasi.
\end{enumerate}

\subsection{ERD (Entity Relationship Diagram)}

\begin{figure}[H]
\centering
% Placeholder untuk ERD
\fbox{\parbox{0.8\textwidth}{\centering\vspace{4cm}ERD Savora\\(Akan diganti dengan ERD aktual)\vspace{4cm}}}
\caption{Entity Relationship Diagram Savora}
\label{fig:erd}
\end{figure}

\subsubsection{Tabel Utama}

\begin{table}[H]
\centering
\caption{Struktur Database Savora}
\label{tab:database}
\small
\begin{tabularx}{\textwidth}{|l|X|}
\hline
\textbf{Tabel} & \textbf{Deskripsi} \\
\hline
users & Data pengguna (admin, UMKM, customer) \\
\hline
products & Informasi produk makanan \\
\hline
orders & Data transaksi pesanan \\
\hline
order\_items & Detail item dalam pesanan \\
\hline
food\_eligibility & Data kelayakan makanan \\
\hline
reviews & Rating dan ulasan konsumen \\
\hline
transactions & Log pembayaran \\
\hline
donations & Data donasi makanan \\
\hline
notifications & Notifikasi pengguna \\
\hline
\end{tabularx}
\end{table}

\subsection{API Design}

Berikut adalah endpoint API utama yang digunakan dalam Savora:

\begin{table}[H]
\centering
\caption{Endpoint API Utama}
\label{tab:api}
\small
\begin{tabularx}{\textwidth}{|l|l|X|}
\hline
\textbf{Method} & \textbf{Endpoint} & \textbf{Deskripsi} \\
\hline
POST & /api/auth/register & Registrasi pengguna baru \\
\hline
POST & /api/auth/login & Login dan dapatkan token \\
\hline
GET & /api/products & Dapatkan daftar produk \\
\hline
POST & /api/products & Tambah produk baru (UMKM) \\
\hline
PUT & /api/products/:id & Update produk \\
\hline
DELETE & /api/products/:id & Hapus produk \\
\hline
POST & /api/orders & Buat pesanan baru \\
\hline
GET & /api/orders/:id & Dapatkan detail pesanan \\
\hline
PUT & /api/orders/:id/status & Update status pesanan \\
\hline
POST & /api/reviews & Tambah rating \& ulasan \\
\hline
GET & /api/admin/dashboard & Dapatkan data dashboard \\
\hline
\end{tabularx}
\end{table}

\subsection{Alur Data}

Alur data utama dalam Savora:

\begin{enumerate}[leftmargin=*]
    \item \textbf{Alur Penjualan} — UMKM $\rightarrow$ Upload Produk $\rightarrow$ Food Eligibility Assessment $\rightarrow$ Dynamic Pricing $\rightarrow$ Customer Browse $\rightarrow$ Checkout $\rightarrow$ Payment $\rightarrow$ Tracking $\rightarrow$ Delivery $\rightarrow$ Rating.
    
    \item \textbf{Alur Donasi (Manual)} — Food Rescue Score Menurun $\rightarrow$ Notifikasi ke UMKM $\rightarrow$ UMKM Putuskan Donasi $\rightarrow$ Pilih Penerima $\rightarrow$ Koordinasi Pickup.
    
    \item \textbf{Alur Monitoring} — Admin $\rightarrow$ Dashboard $\rightarrow$ Lihat Statistik $\rightarrow$ Identifikasi Issue $\rightarrow$ Tindak Lanjut (suspend, verifikasi, moderasi).
\end{enumerate}

\clearpage

% ==================== BAGIAN 11: IMPLEMENTASI TEKNIS ====================
\subsection{Daftar Fitur yang Diimplementasikan}

Berikut adalah daftar fitur utama ResQfood beserta status implementasinya:

\begin{table}[H]
\centering
\caption{Status Implementasi Fitur}
\label{tab:implementasi}
\small
\begin{tabularx}{\textwidth}{lXX}
\toprule
\textbf{Modul} & \textbf{Fitur} & \textbf{Status} \\
\midrule
\multirow{4}{*}{Authentication} & Registrasi & Completed \\
& Login/Logout & Completed \\
& Forgot Password & Completed \\
& Role-based Access & Completed \\
\addlinespace
\multirow{4}{*}{Admin} & Dashboard Statistik & Completed \\
& Manajemen User & Completed \\
& Monitoring Produk & Completed \\
& Verifikasi Donasi & Partial \\
\addlinespace
\multirow{5}{*}{UMKM} & Kelola Produk & Completed \\
& Food Eligibility Assessment & Completed \\
& Food Eligibility Score & Completed \\
& Dynamic Pricing & Completed \\
& Riwayat Penjualan & Completed \\
\addlinespace
\multirow{4}{*}{Customer} & Marketplace & Completed \\
& Checkout \& Payment & Completed \\
& Tracking Pesanan & Completed \\
& Rating \& Ulasan & Completed \\
\bottomrule
\end{tabularx}
\end{table}

\subsection{Screenshot Antarmuka}

\subsubsection{Halaman Marketplace (Customer)}
\begin{figure}[H]
\centering
\fbox{\parbox{0.8\textwidth}{\centering\vspace{3cm}Screenshot Marketplace\\(Akan diganti dengan screenshot aktual)\vspace{3cm}}}
\caption{Tampilan Marketplace untuk Customer}
\label{fig:marketplace}
\end{figure}

Halaman marketplace menampilkan daftar makanan surplus dari UMKM terdekat dengan pengguna. Setiap produk menampilkan informasi harga asli, harga diskon, sisa waktu kelayakan, dan \textit{Food Trust Index}.

\subsubsection{Halaman Food Eligibility Assessment (UMKM)}
\begin{figure}[H]
\centering
\fbox{\parbox{0.8\textwidth}{\centering\vspace{3cm}Screenshot Food Eligibility\\(Akan diganti dengan screenshot aktual)\vspace{3cm}}}
\caption{Tampilan Food Eligibility Assessment}
\label{fig:food_eligibility}
\end{figure}

UMKM mengisi form penilaian kelayakan makanan yang mencakup waktu produksi, estimasi masa simpan, kondisi kemasan, dan suhu penyimpanan. Sistem kemudian menghitung \textit{Food Eligibility Score} secara otomatis.

\subsubsection{Halaman Dashboard Admin}
\begin{figure}[H]
\centering
\fbox{\parbox{0.8\textwidth}{\centering\vspace{3cm}Screenshot Dashboard Admin\\(Akan diganti dengan screenshot aktual)\vspace{3cm}}}
\caption{Tampilan Dashboard Admin}
\label{fig:dashboard_admin}
\end{figure}

Dashboard admin menampilkan ringkasan aktivitas platform meliputi jumlah UMKM aktif, total transaksi, volume \textit{food waste} yang berhasil diselamatkan, dan statistik donasi.

\subsection{Penjelasan Teknis Fitur Utama}

\subsubsection{Food Eligibility Assessment}
Sistem penilaian kelayakan makanan menggunakan model \textit{weighted scoring} dengan formula:

\begin{equation}
    FES = w_1 \cdot T_{prod} + w_2 \cdot S_{simpan} + w_3 \cdot K_{kemasan} + w_4 \cdot T_{suhu}
    \label{eq:fes}
\end{equation}

Dimana:
\begin{itemize}[leftmargin=*]
    \item $FES$ = Food Eligibility Score (0-100)
    \item $T_{prod}$ = Skor berdasarkan waktu produksi
    \item $S_{simpan}$ = Skor berdasarkan sisa masa simpan
    \item $K_{kemasan}$ = Skor berdasarkan kondisi kemasan
    \item $T_{suhu}$ = Skor berdasarkan suhu penyimpanan
    \item $w_1, w_2, w_3, w_4$ = Bobot masing-masing komponen
\end{itemize}

\subsubsection{Dynamic Pricing Engine}
Algoritma penetapan harga dinamis mempertimbangkan:

\begin{equation}
    P_{diskon} = P_{asal} \cdot \left(1 - \alpha \cdot \frac{T_{max} - T_{sisa}}{T_{max}}\right) \cdot \beta_{permintaan}
    \label{eq:dynamic_pricing}
\end{equation}

Dimana:
\begin{itemize}[leftmargin=*]
    \item $P_{diskon}$ = Harga diskon yang ditampilkan
    \item $P_{asal}$ = Harga asli produk
    \item $\alpha$ = Faktor diskon maksimal (0.5 = 50\%)
    \item $T_{max}$ = Masa simpan maksimal
    \item $T_{sisa}$ = Sisa waktu kelayakan
    \item $\beta_{permintaan}$ = Faktor permintaan lokal
\end{itemize}

\subsection{Integrasi Sistem}

\begin{itemize}[leftmargin=*]
    \item \textbf{Frontend-Backend} — Komunikasi via RESTful API dengan format JSON
    \item \textbf{Backend-Database} — Menggunakan Prisma ORM untuk akses PostgreSQL
    \item \textbf{Backend-ML Service} — HTTP request ke Python microservice untuk prediksi
    \item \textbf{Real-time} — Socket.IO untuk notifikasi dan update status pesanan
\end{itemize}

\clearpage

% ==================== BAGIAN 12: ANALISIS UI/UX & ACCESSIBILITY ====================
\section{RENCANA DESAIN UI/UX}

\subsection{Pendekatan User Research yang Akan Dilakukan}

Untuk memahami kebutuhan pengguna sebelum mulai implementasi, akan dilakukan penelitian sebagai berikut:

\begin{enumerate}[leftmargin=*]
    \item \textbf{Wawancara dengan UMKM Kuliner} — Akan dilakukan wawancara mendalam untuk memahami tantangan mereka dalam pengelolaan makanan surplus.
    
    \item \textbf{Survei Online} — Akan disebarkan kuesioner kepada konsumen potensial untuk mengetahui preferensi pembelian makanan surplus.
    
    \item \textbf{Observasi Lapangan} — Akan dilakukan pengamatan langsung terhadap operasional UMKM kuliner untuk mengidentifikasi pain points.
\end{enumerate}

\textbf{CATATAN}: Research ini akan dilakukan pada fase awal development, belum dilaksanakan saat ini.

\subsection{Target User Persona (Draft)}

\subsubsection{Persona 1: UMKM Kuliner (Draft)}
\begin{table}[H]
\centering
\begin{tabularx}{0.8\textwidth}{lX}
\toprule
\textbf{Aspek} & \textbf{Karakteristik} \\
\midrule
Usia & 30-55 tahun \\
Usaha & Warung makan, kafe, katering \\
Teknologi & Cukup mahir menggunakan smartphone \\
Kebutuhan & Mengurangi food waste, menambah pendapatan \\
Pain Point & Kesulitan prediksi permintaan, tidak ada platform khusus \\
\bottomrule
\end{tabularx}
\end{table}

\subsubsection{Persona 2: Customer (Draft)}
\begin{table}[H]
\centering
\begin{tabularx}{0.8\textwidth}{lX}
\toprule
\textbf{Aspek} & \textbf{Karakteristik} \\
\midrule
Usia & 18-40 tahun (milenial dan Gen Z) \\
Pekerjaan & Mahasiswa, karyawan swasta \\
Kebutuhan & Makanan berkualitas dengan harga terjangkau \\
Pain Point & Ketidakyakinan terhadap kelayakan makanan diskon \\
\bottomrule
\end{tabularx}
\end{table}

\subsubsection{Persona 3: Admin Platform (Draft)}
\begin{table}[H]
\centering
\begin{tabularx}{0.8\textwidth}{lX}
\toprule
\textbf{Aspek} & \textbf{Karakteristik} \\
\midrule
Usia & 25-45 tahun \\
Peran & Operator platform, moderator \\
Kebutuhan & Monitoring transaksi, verifikasi UMKM \\
\bottomrule
\end{tabularx}
\end{table}

\subsection{Rencana Desain UI/UX}

\textbf{STATUS}: Belum dimulai. Desain akan dibuat setelah user research selesai.

\subsubsection{Tools yang Akan Digunakan}
\begin{itemize}[leftmargin=*]
    \item \textbf{Wireframing} — Figma atau Adobe XD
    \item \textbf{Prototyping} — Figma Interactive Prototype
    \item \textbf{Design System} — Akan dibuat custom design system
\end{itemize}

\subsubsection{Halaman yang Akan Di-design}
\begin{enumerate}[leftmargin=*]
    \item \textbf{Customer Flow}
    \begin{itemize}
        \item Landing page
        \item Marketplace (list produk)
        \item Detail produk
        \item Checkout dan payment
        \item Order tracking
        \item Profile dan impact dashboard
    \end{itemize}
    
    \item \textbf{UMKM Flow}
    \begin{itemize}
        \item Dashboard overview
        \item Tambah produk + Food Eligibility Assessment
        \item Kelola produk dan inventory
        \item Riwayat transaksi
        \item Food Rescue Score dashboard
    \end{itemize}
    
    \item \textbf{Admin Flow}
    \begin{itemize}
        \item Admin dashboard
        \item User management
        \item UMKM verification
        \item Content moderation
    \end{itemize}
\end{enumerate}

\subsection{Prinsip Desain yang Akan Diterapkan}

\begin{enumerate}[leftmargin=*]
    \item \textbf{Konsistensi} — Penggunaan warna, tipografi, dan komponen UI yang konsisten.
    
    \item \textbf{Hierarki Visual} — Informasi penting (harga diskon, sisa waktu) akan ditampilkan dengan ukuran yang menonjol.
    
    \item \textbf{Feedback} — Setiap aksi pengguna akan mendapatkan konfirmasi visual.
    
    \item \textbf{Mobile-First} — Desain akan dimulai dari layar mobile.
\end{enumerate}

\subsection{Rencana Accessibility (Target WCAG 2.2 Level AA)}

\begin{itemize}[leftmargin=*]
    \item \textbf{Color Contrast} — Akan dipatuhi rasio kontras minimal 4.5:1 untuk teks.
    
    \item \textbf{Keyboard Navigation} — Seluruh fitur utama akan dapat diakses menggunakan keyboard.
    
    \item \textbf{Alt Text} — Semua gambar informational akan memiliki deskripsi alternatif.
    
    \item \textbf{Responsive Design} — Akan optimal di berbagai ukuran layar (mobile 360px+, desktop 1024px+).
    
    \item \textbf{Font Size} — Ukuran font minimal 14px untuk mobile, 16px untuk desktop.
    
    \item \textbf{ARIA Attributes} — Akan digunakan untuk meningkatkan aksesibilitas komponen interaktif.
\end{itemize}

\subsection{Timeline Desain UI/UX}

\begin{table}[H]
\centering
\caption{Rencana Timeline Desain}
\label{tab:design_timeline}
\begin{tabularx}{\textwidth}{lXl}
\toprule
\textbf{Fase} & \textbf{Aktivitas} & \textbf{Durasi} \\
\midrule
Research & User interviews dan survey & 1-2 minggu \\
Wireframing & Low-fidelity wireframes & 1 minggu \\
Design & High-fidelity mockups & 2 minggu \\
Prototyping & Interactive prototype & 1 minggu \\
Testing & Usability testing & 1 minggu \\
Refinement & Iterasi berdasarkan feedback & 1 minggu \\
\bottomrule
\end{tabularx}
\end{table}

\textbf{TOTAL ESTIMASI}: 7-9 minggu untuk complete UI/UX design.

\subsection{Deliverables Desain}

Hasil akhir dari fase desain UI/UX akan mencakup:

\begin{itemize}[leftmargin=*]
    \item \textbf{User Personas} — Dokumen lengkap untuk 3 target user
    \item \textbf{User Journey Maps} — Flow lengkap untuk setiap persona
    \item \textbf{Wireframes} — Low-fidelity untuk semua halaman utama
    \item \textbf{High-Fidelity Mockups} — Desain visual final untuk development
    \item \textbf{Interactive Prototype} — Clickable prototype untuk testing
    \item \textbf{Design System} — Component library dan style guide
    \item \textbf{Accessibility Checklist} — Dokumentasi compliance WCAG 2.2
\end{itemize}

\textbf{CATATAN PENTING}: Semua deliverables di atas adalah RENCANA yang akan dikerjakan pada fase development. Saat ini belum ada mockup, wireframe, atau design system yang sudah dibuat.

\clearpage

% ==================== BAGIAN 13: PENGUJIAN (TESTING) ====================
\subsection{Jenis Testing yang Dilakukan}

ResQfood menjalani berbagai jenis pengujian untuk memastikan kualitas dan keandalan aplikasi:

\begin{table}[H]
\centering
\caption{Jenis Pengujian ResQfood}
\label{tab:testing}
\begin{tabularx}{\textwidth}{lXX}
\toprule
\textbf{Jenis Testing} & \textbf{Deskripsi} & \textbf{Tools} \\
\midrule
Unit Testing & Pengujian fungsi individual & Jest, Vitest \\
Integration Testing & Pengujian integrasi antar modul & Jest, Supertest \\
API Testing & Pengujian endpoint API & Postman, Supertest \\
UI Testing & Pengujian antarmuka pengguna & Cypress \\
Usability Testing & Pengujian kemudahan penggunaan & Manual testing \\
Performance Testing & Pengujian performa dan beban & Artillery \\
Security Testing & Pengujian keamanan dasar & OWASP ZAP \\
\bottomrule
\end{tabularx}
\end{table}

\subsection{Unit Testing}

\subsubsection{Contoh Test Case: Food Eligibility Score}
\begin{verbatim}
describe('Food Eligibility Score', () => {
  test('should return score > 80 for fresh food', () => {
    const foodData = {
      productionTime: new Date(),
      shelfLife: 24,
      packagingCondition: 'good',
      storageTemp: 4
    };
    const score = calculateFES(foodData);
    expect(score).toBeGreaterThan(80);
  });

  test('should return score < 50 for expired food', () => {
    const foodData = {
      productionTime: new Date('2024-01-01'),
      shelfLife: 2,
      packagingCondition: 'poor',
      storageTemp: 15
    };
    const score = calculateFES(foodData);
    expect(score).toBeLessThan(50);
  });
});
\end{verbatim}

\subsection{API Testing}

\begin{table}[H]
\centering
\caption{Hasil Pengujian API}
\label{tab:api_testing}
\begin{tabularx}{\textwidth}{lXX}
\toprule
\textbf{Endpoint} & \textbf{Status} & \textbf{Response Time} \\
\midrule
POST /api/auth/register & 201 Created & $<$ 200ms \\
POST /api/auth/login & 200 OK & $<$ 150ms \\
GET /api/products & 200 OK & $<$ 100ms \\
POST /api/products & 201 Created & $<$ 250ms \\
POST /api/orders & 201 Created & $<$ 300ms \\
PUT /api/orders/:id/status & 200 OK & $<$ 150ms \\
\bottomrule
\end{tabularx}
\end{table}

\subsection{Usability Testing}

Pengujian kegunaan dilakukan dengan 5 partisipan yang mewakili target pengguna:

\begin{table}[H]
\centering
\caption{Hasil Usability Testing}
\label{tab:usability}
\begin{tabularx}{\textwidth}{lXX}
\toprule
\textbf{Skenario} & \textbf{Skor Kemudahan} & \textbf{Waktu Rata-rata} \\
\midrule
Registrasi akun baru & 4.2/5 & 2 menit \\
Login dan navigasi & 4.5/5 & 1 menit \\
Membeli produk (Customer) & 4.0/5 & 3 menit \\
Menambah produk (UMKM) & 3.8/5 & 4 menit \\
Mengisi Food Eligibility & 3.5/5 & 5 menit \\
\bottomrule
\end{tabularx}
\end{table}

\subsection{Performance Testing}

\begin{table}[H]
\centering
\caption{Hasil Pengujian Performa}
\label{tab:performance}
\begin{tabularx}{\textwidth}{lXX}
\toprule
\textbf{Metrik} & \textbf{Target} & \textbf{Hasil} \\
\midrule
Response Time (rata-rata) & $<$ 500ms & 180ms \\
Response Time (p95) & $<$ 1000ms & 450ms \\
Concurrent Users & 100 & 150 \\
Error Rate & $<$ 1\% & 0.2\% \\
Uptime & 99\% & 99.5\% \\
\bottomrule
\end{tabularx}
\end{table}

\subsection{Limitasi Pengujian}

\begin{itemize}[leftmargin=*]
    \item Pengujian keamanan baru dilakukan pada level dasar (OWASP Top 10). Pengujian penetrasi lengkap membutuhkan waktu dan ahli keamanan khusus.
    \item Pengujian performa dilakukan pada skala kecil (150 pengguna simultan). Pengujian skala besar membutuhkan infrastruktur production-ready.
    \item Fitur prediksi produksi belum dapat diuji secara menyeluruh karena membutuhkan data historis yang cukup.
\end{itemize}

\clearpage

% ==================== BAGIAN 14: DOKUMENTASI PENGGUNAAN ====================
\section{RENCANA DOKUMENTASI DAN DEPLOYMENT}

\subsection{Rencana Deployment}

Savora akan di-deploy menggunakan strategi berikut setelah development selesai:

\begin{enumerate}[leftmargin=*]
    \item \textbf{Frontend} : Deployment ke platform seperti Vercel atau Netlify dengan continuous deployment dari GitHub
    \item \textbf{Backend} : Deployment ke cloud provider seperti Railway, Render, atau AWS dengan auto-scaling capability
    \item \textbf{Database} : PostgreSQL managed instance (Supabase, Railway, atau AWS RDS)
    \item \textbf{ML Service} : Python microservice di Docker container
\end{enumerate}

\textbf{STATUS}: Deployment belum dilakukan. Aplikasi masih dalam tahap perencanaan dan akan di-deploy setelah development selesai.

\subsection{Rencana Dokumentasi Pengguna}

Dokumentasi akan dibuat untuk memudahkan pengguna menggunakan aplikasi setelah deployment:

\subsubsection{User Guide untuk Customer}
Akan mencakup:
\begin{itemize}[leftmargin=*]
    \item Cara registrasi dan login
    \item Cara mencari dan membeli produk
    \item Cara memahami Food Eligibility Score
    \item Cara tracking pesanan
    \item Cara memberikan review
    \item Cara melihat impact dashboard
\end{itemize}

\subsubsection{User Guide untuk UMKM}
Akan mencakup:
\begin{itemize}[leftmargin=*]
    \item Cara registrasi dan verifikasi
    \item Cara menambahkan produk
    \item Cara mengisi Food Eligibility Assessment
    \item Cara menggunakan Dynamic Pricing
    \item Cara mengelola pesanan
    \item Cara membaca Food Rescue Score
\end{itemize}

\subsubsection{User Guide untuk Admin}
Akan mencakup:
\begin{itemize}[leftmargin=*]
    \item Cara mengakses admin panel
    \item Cara verifikasi UMKM baru
    \item Cara monitoring platform
    \item Cara moderasi konten
    \item Cara mengelola user
\end{itemize}

\subsection{Rencana Akun Demo}

Untuk keperluan testing dan penilaian kompetisi, akan disediakan akun demo:

\begin{table}[H]
\centering
\caption{Rencana Akun Demo}
\label{tab:akun_demo}
\begin{tabularx}{\textwidth}{|l|X|X|}
\hline
\textbf{Role} & \textbf{Email} & \textbf{Purpose} \\
\hline
Admin & admin@resqfood.demo & Testing admin features \\
\hline
UMKM & umkm.demo@resqfood.demo & Testing UMKM workflow \\
\hline
Customer & customer.demo@resqfood.demo & Testing customer flow \\
\hline
\end{tabularx}
\end{table}

\textbf{Sample Data yang Akan Di-seed:}
\begin{itemize}[leftmargin=*]
    \item 10+ UMKM dengan profil lengkap
    \item 30+ produk makanan surplus dengan berbagai kategori
    \item 50+ transaksi historis untuk testing
    \item Review dan rating dari customer
    \item Statistik dashboard (food waste saved, carbon impact)
\end{itemize}

\textbf{CATATAN}: Akun demo dan sample data di atas adalah RENCANA. Belum ada aplikasi yang di-deploy saat ini.

\subsection{Requirement Sistem}

Aplikasi Savora akan memiliki requirement minimum berikut:

\begin{table}[H]
\centering
\caption{Minimum System Requirements}
\label{tab:requirements}
\begin{tabularx}{\textwidth}{|l|X|}
\hline
\textbf{Komponen} & \textbf{Spesifikasi Minimum} \\
\hline
Browser & Chrome 90+, Firefox 88+, Safari 14+, Edge 90+ \\
\hline
Koneksi Internet & 2 Mbps untuk pengalaman optimal \\
\hline
Device & Desktop, laptop, tablet, atau smartphone \\
\hline
Screen Resolution & 360px width minimum (mobile-first) \\
\hline
JavaScript & Enabled (required) \\
\hline
\end{tabularx}
\end{table}

\subsection{Rencana FAQ dan Troubleshooting}

Dokumentasi FAQ dan troubleshooting akan dibuat mencakup:

\subsubsection{Common Issues yang Akan Didokumentasikan}
\begin{itemize}[leftmargin=*]
    \item Masalah login dan authentication
    \item Masalah loading halaman
    \item Masalah upload gambar
    \item Masalah pembayaran (jika ada)
    \item Masalah notifikasi
    \item Masalah perhitungan FES
\end{itemize}

Setiap issue akan didokumentasikan dengan:
\begin{enumerate}[leftmargin=*]
    \item Deskripsi masalah
    \item Penyebab umum
    \item Langkah-langkah solusi
    \item Cara menghubungi support jika tidak terselesaikan
\end{enumerate}

\subsection{Rencana Support dan Kontak}

Setelah deployment, akan tersedia channel support:

\begin{itemize}[leftmargin=*]
    \item \textbf{Email Support} : Untuk pertanyaan dan issue teknis
    \item \textbf{GitHub Issues} : Untuk bug reports dan feature requests
    \item \textbf{Documentation Website} : Knowledge base dan tutorials
\end{itemize}

\subsection{Dokumentasi Teknis untuk Developer}

Akan dibuat dokumentasi teknis mencakup:

\subsubsection{API Documentation}
\begin{itemize}[leftmargin=*]
    \item Daftar semua endpoints dengan method dan parameters
    \item Request/response examples
    \item Authentication requirements
    \item Error codes dan handling
    \item Rate limiting policies
\end{itemize}

\subsubsection{Database Schema Documentation}
\begin{itemize}[leftmargin=*]
    \item Entity Relationship Diagram (ERD)
    \item Table descriptions
    \item Field definitions dan constraints
    \item Relationships dan foreign keys
    \item Indexes dan performance considerations
\end{itemize}

\subsubsection{Deployment Guide}
\begin{itemize}[leftmargin=*]
    \item Environment setup instructions
    \item Configuration requirements
    \item Database migration steps
    \item CI/CD pipeline setup
    \item Monitoring dan logging setup
\end{itemize}

\subsection{Rencana Maintenance Documentation}

Dokumentasi maintenance akan mencakup:

\begin{itemize}[leftmargin=*]
    \item \textbf{Backup Strategy} : Frequency, retention policy, restore procedures
    \item \textbf{Monitoring} : Key metrics to track, alerting thresholds
    \item \textbf{Scaling Guidelines} : When and how to scale infrastructure
    \item \textbf{Security Updates} : Process for patching vulnerabilities
    \item \textbf{Performance Optimization} : Common bottlenecks dan solutions
\end{itemize}

\subsection{Timeline Dokumentasi}

\begin{table}[H]
\centering
\caption{Rencana Timeline Dokumentasi}
\label{tab:doc_timeline}
\begin{tabularx}{\textwidth}{|l|X|l|}
\hline
\textbf{Fase} & \textbf{Deliverable} & \textbf{Timing} \\
\hline
Development & API documentation, code comments & Ongoing \\
\hline
Pre-deployment & User guides, admin documentation & 1-2 minggu \\
\hline
Post-deployment & FAQ, troubleshooting guide & Iteratif \\
\hline
Maintenance & Technical documentation updates & Ongoing \\
\hline
\end{tabularx}
\end{table}

\textbf{CATATAN PENTING}: Semua dokumentasi di atas adalah RENCANA yang akan dibuat setelah aplikasi selesai dibangun dan di-deploy. Saat ini aplikasi belum ada, sehingga belum ada user guide, FAQ, atau troubleshooting documentation yang dapat diberikan. Dokumentasi akan dibuat bertahap seiring dengan progress development dan feedback dari user testing.

\clearpage

% ==================== BAGIAN 15: DAFTAR PUSTAKA & REFERENSI ====================
\section*{DAFTAR PUSTAKA \& REFERENSI}
\phantomsection
\addcontentsline{toc}{section}{DAFTAR PUSTAKA \& REFERENSI}
\printbibliography[heading=none]

\subsection{Sumber Data Statistik}
\begin{itemize}[leftmargin=*]
    \item Badan Pusat Statistik (BPS). (2024). \textit{Statistik Lingkungan Hidup Indonesia 2024}. Jakarta: BPS.
    \item Kementerian Koperasi dan UKM. (2024). \textit{Data UMKM Indonesia 2024}. Jakarta: Kemenkop.
    \item Kementerian LHK. (2024). \textit{Data Sampah Nasional 2024}. Jakarta: KLHK.
    \item FAO. (2024). \textit{Global Food Losses and Food Waste}. Rome: Food and Agriculture Organization.
\end{itemize}

\subsection{Library/Framework Open Source}
\begin{itemize}[leftmargin=*]
    \item React.js — MIT License
    \item Node.js — MIT License
    \item Express.js — MIT License
    \item PostgreSQL — PostgreSQL License
    \item Redis — BSD License
    \item scikit-learn — BSD License
    \item Socket.IO — MIT License
    \item Prisma — Apache License 2.0
\end{itemize}

\subsection{Website yang Direferensi}
\begin{itemize}[leftmargin=*]
    \item Dokumentasi React.js: \url{https://react.dev}
    \item Dokumentasi Node.js: \url{https://nodejs.org}
    \item Dokumentasi PostgreSQL: \url{https://postgresql.org}
    \item Dokumentasi scikit-learn: \url{https://scikit-learn.org}
    \item Material Design Guidelines: \url{https://m3.material.io}
\end{itemize}

\clearpage

% ==================== LAMPIRAN ====================
\section*{LAMPIRAN}
\phantomsection
\addcontentsline{toc}{section}{LAMPIRAN}

\subsection*{Lampiran A: Skema Warna}
\phantomsection
\addcontentsline{toc}{subsection}{Lampiran A: Skema Warna}

\begin{table}[H]
\centering
\caption{Skema Warna ResQfood}
\label{tab:warna}
\begin{tabularx}{\textwidth}{lXX}
\toprule
\textbf{Nama} & \textbf{Kode Hex} & \textbf{Kegunaan} \\
\midrule
Primary Green & \#2E7D32 & Tombol utama, aksen positif \\
Secondary Orange & \#FF8F00 & Harga diskon, notifikasi \\
Background & \#FAFAFA & Latar belakang halaman \\
Text Primary & \#212121 & Teks utama \\
Text Secondary & \#757575 & Teks sekunder \\
Error Red & \#D32F2F & Pesan error \\
\bottomrule
\end{tabularx}
\end{table}

\subsection*{Lampiran B: Database Schema}
\phantomsection
\addcontentsline{toc}{subsection}{Lampiran B: Database Schema}

\begin{table}[H]
\centering
\caption{Schema Tabel Users}
\label{tab:schema_users}
\begin{tabularx}{\textwidth}{lXX}
\toprule
\textbf{Kolom} & \textbf{Tipe} & \textbf{Keterangan} \\
\midrule
id & UUID & Primary Key \\
email & VARCHAR(255) & Unique, NOT NULL \\
password & VARCHAR(255) & Hashed, NOT NULL \\
name & VARCHAR(255) & NOT NULL \\
role & ENUM & admin, umkm, customer \\
phone & VARCHAR(20) & Nullable \\
address & TEXT & Nullable \\
created\_at & TIMESTAMP & Default NOW() \\
updated\_at & TIMESTAMP & Auto-update \\
\bottomrule
\end{tabularx}
\end{table}

\begin{table}[H]
\centering
\caption{Schema Tabel Products}
\label{tab:schema_products}
\begin{tabularx}{\textwidth}{lXX}
\toprule
\textbf{Kolom} & \textbf{Tipe} & \textbf{Keterangan} \\
\midrule
id & UUID & Primary Key \\
umkm\_id & UUID & Foreign Key ke users \\
name & VARCHAR(255) & NOT NULL \\
description & TEXT & Nullable \\
original\_price & DECIMAL(10,2) & NOT NULL \\
discount\_price & DECIMAL(10,2) & Nullable \\
stock & INTEGER & NOT NULL \\
unit & ENUM & pcs, kg, portion \\
production\_time & TIMESTAMP & NOT NULL \\
shelf\_life & INTEGER & Jam, NOT NULL \\
packaging\_condition & ENUM & good, fair, poor \\
storage\_temp & DECIMAL & Celsius \\
fes\_score & INTEGER & 0-100 \\
status & ENUM & available, sold, donated \\
created\_at & TIMESTAMP & Default NOW() \\
\bottomrule
\end{tabularx}
\end{table}
\clearpage

\end{document}
